% ==========================
% Introduction
% ==========================
\clearpage
\newgeometry{
	top=2.1cm,
	bottom=5.4cm,
	inner=4.6cm,
	outer=2.3cm,
	headsep=0.8cm,
	includehead,
	twoside
}
\reversemarginpar
\setlength{\headwidth}{\textwidth}
\fancyhfoffset[LE,RO]{0pt}
\pagenumbering{arabic}
\setcounter{page}{1}
\pagestyle{main}
\CustomIntroduction{Introduction}
\label{ch:intro}

\section*{General Context}

General relativity, formulated by Einstein in 1915, remains the cornerstone theory describing gravitation as the curvature of spacetime induced by the presence of mass and energy \cite{einstein1916}. Since its formulation, it has been confirmed by numerous experimental and observational tests, ranging from the perihelion precession of Mercury to the direct detection of gravitational waves \cite{ligo2016}.

The study of exact solutions to the Einstein field equations, such as the Schwarzschild and Kerr metrics, continues to play a central role in understanding compact astrophysical objects such as black holes and neutron stars \cite{misner1973}.

\section*{Problem Statement}

Despite a century of theoretical and observational progress, several open questions remain, including the behavior of spacetime near singularities, the compatibility of general relativity with quantum mechanics, and the interpretation of dark energy within the cosmological constant framework \cite{wald1984}. This work addresses some of these questions by proposing a detailed study of specific exact solutions of the Einstein field equations and their physical implications.

\section*{Objectives}

The main objective of this thesis is to analyze, interpret, and numerically explore selected exact solutions of general relativity. Secondary objectives include studying the geodesic structure associated with these solutions and comparing their predictions with astrophysical observations.

\begin{center}
	\ObjectiveBox{%
		The main objective of this work is to conduct a detailed theoretical and numerical study of exact solutions of the Einstein field equations, with particular emphasis on the Schwarzschild and Kerr metrics, in order to analyze their geodesic structure, horizon properties, and physical consequences for the behavior of matter and light in strong gravitational fields.
	}
\end{center}

\section*{Scientific Context}

The rapid development of relativistic astrophysics, particularly following the direct detection of gravitational waves and the imaging of black hole shadows, has made the study of strong-field gravity a highly active research area. From compact binary mergers to the dynamics of matter near event horizons, general relativity provides the essential theoretical framework for interpreting these phenomena.

Exact solutions of the Einstein field equations have emerged as indispensable tools for modeling astrophysical compact objects, owing to their ability to capture the nonlinear geometric structure of spacetime in regions of strong gravity. Their confirmation through observations such as the Event Horizon Telescope imaging of M87* has established them as central benchmarks of modern relativistic astrophysics \cite{eht2019}.

% ==========================
% Last Line Shape 
% ==========================
\perfectdivider
